\section{Machine Learning in Astrophysics}
\label{sec:3.3}

The SBI framework described in Section~\ref{sec:3.2} already uses machine learning to approximate intractable posteriors, but machine learning plays a broader role in this thesis.
In Chapter~\ref{ch:6}, a convolutional neural network reconstructs the source-count distribution directly from gamma-ray sky maps; in Chapter~\ref{ch:5}, density estimation and mixture modeling are used to search for a dark matter subhalo population among unassociated Fermi-LAT sources.
This section develops the tools needed for both analyses: the relevant learning tasks (Section~\ref{sec:3.3.1}), neural network architectures and uncertainty estimation (Section~\ref{sec:3.3.2}), and non-parametric density estimation (Section~\ref{sec:3.3.3}).

\subsection{Learning Tasks}
\label{sec:3.3.1}

Machine learning provides a flexible framework for extracting information from data when explicit physical models are either unavailable or computationally intractable \cite{bishop,Hastie:2009}.
The problems addressed in this thesis can be organized into three canonical learning tasks, each motivated by a distinct physical question.

The first is \textit{classification}: assigning a discrete label to an input based on its observed features.
In the gamma-ray context, this amounts to determining the astrophysical nature of a detected source.
Roughly one-third of the sources in Fermi-LAT catalogs lack firm multi-wavelength associations \cite{Bhat:2022}, and supervised classification algorithms---trained on populations of identified AGNs and pulsars---are used to predict the class membership of these unassociated objects.
The input features are quantities derived from the gamma-ray signal itself: spectral curvature significance, variability index, hardness ratios, and energy flux \cite{Bhat:2022}.
In Chapter~\ref{ch:5}, this paradigm is extended beyond standard source classes through a generative mixture model that accommodates an entirely new population---dark matter subhalos---among the unassociated sources.

The second task is \textit{regression}: predicting a continuous quantity from observational data.
In Chapter~\ref{ch:6}, a convolutional neural network estimates the source-count distribution $dN/dS$ directly from gamma-ray sky maps.
By adopting a heteroscedastic loss function (described in Section~\ref{sec:3.3.2}), the network simultaneously performs regression on both the $dN/dS$ values and their associated uncertainties, predicting the variance of each output alongside its mean \cite{heteroscedastic-error}.

The third task is \textit{density estimation}: approximating the probability distribution from which observed source features are drawn.
While classification asks ``what class is this source?
'', density estimation asks the complementary question ``what does the overall population look like?'' This shift from individual classification to population-level inference underpins the generative mixture model developed in Chapter~\ref{ch:5}, where Kernel Density Estimation---used to model the feature distributions of associated sources---is combined with mixture model fitting to infer the energy spectrum composition of the unassociated gamma-ray sources.



% where Kernel Density Estimation is used to model the spectral parameter distributions of known source classes.
% The resulting class-conditional densities $p(\mathbf{x}|k)$ are then combined into a generative mixture model whose parameters---including the prevalence of a potential dark matter component---are determined by maximizing the likelihood over the unassociated source data.
% The mathematical framework for density estimation and mixture models is developed in Section~\ref{sec:3.3.3}.

\subsection{Neural Networks and Deep Learning}
\label{sec:3.3.2}

A neural network is a parametric function $f_{\boldsymbol{w}}: \mathbb{R}^n \to \mathbb{R}^m$ composed of a sequence of affine transformations, each followed by a non-linear activation function.
In a network with $L$ layers, the output of layer $\ell$ is $\mathbf{h}^{(\ell)} = \sigma(\mathbf{W}^{(\ell)}\mathbf{h}^{(\ell-1)} + \mathbf{b}^{(\ell)})$, where $\mathbf{W}^{(\ell)}$ and $\mathbf{b}^{(\ell)}$ are trainable weight matrices and bias vectors, and $\sigma$ is a non-linear function such as the rectified linear unit $\mathrm{ReLU}(x) = \max(0, x)$.
Given a sufficiently large number of parameters, neural networks can approximate any continuous function to arbitrary accuracy \cite{bishop}, and training consists of adjusting the weights $\boldsymbol{w}$ to minimize a cost function $\mathcal{J}(\boldsymbol{w})$ using gradient-based optimization.

\paragraph{Convolutional neural networks.}

When the input is spatially structured---as is the case for images and sky maps---convolutional neural networks (CNNs) exploit this structure through two properties: weight sharing and translation equivariance.
A convolutional layer applies a set of learnable filters to local patches of the input, producing feature maps that respond to spatial patterns at specific scales.
Because the same filter is applied everywhere across the input, the number of parameters per layer is independent of the input size, and the response to a pattern is the same regardless of its location.
This makes CNNs particularly well-suited to gamma-ray astronomy, where point sources are sparse, approximately isotropic across the sky, and convolved with the instrument's point spread function.

A specific challenge arises when the data live on the sphere rather than on a flat grid.
Three broad strategies have been developed to extend convolutional architectures to spherical data.

The first strategy performs convolutions in the spherical harmonic domain.
Because the spherical harmonics are the natural eigenbasis of the Laplace--Beltrami operator on $S^2$, a spectral convolution can be defined in exact analogy with the classical convolution theorem on the real line.
\texttt{DeepSphere} \cite{DBLP:journals/corr/abs-2012-15000} implements this idea using a graph-based approximation, representing the sphere as a graph whose nodes are HEALPix pixels and whose edges encode pixel adjacencies; spectral graph convolutions then replace the standard spatial filters.
This approach respects the global topology of the sphere and naturally handles irregular pixelizations.

The second strategy applies convolutions to local projections of the sphere.
Rather than operating on the curved manifold directly, these methods project spherical patches onto tangent planes or polyhedral faces, where standard 2D convolutions apply without modification.
Icosahedral projections \cite{eder2019convolutions} decompose the sphere into 20 equilateral triangular faces that cover it nearly uniformly, limiting the distortion inherent in any flat projection.

The third strategy, and the one most directly relevant to this thesis, operates on the HEALPix pixelization itself.
\texttt{NNHealpix} \cite{Krachmalnicoff:2019} exploits the fixed neighborhood structure of HEALPix pixels by unrolling each pixel and its eight nearest neighbors into a nine-element vector and applying standard 1D convolutions with stride nine.
This preserves the full topology of the sphere at the cost of higher computational overhead and sensitivity to the orientation of pixel neighborhoods.
The \texttt{map2patch} algorithm developed in Chapter~\ref{ch:6} addresses these limitations by subdividing the sphere into the 12 equal-area base pixels of the HEALPix scheme and treating each as an independent flat image processed through standard 2D convolutions.
This subdivision preserves the exact area and photon content of every pixel without resampling, at the cost of discarding spatial information at patch boundaries.
For statistically isotropic, small-scale signals such as the point-source population studied in Chapter~\ref{ch:6}, \texttt{map2patch} and \texttt{NNHealpix} produce compatible results, while \texttt{map2patch} achieves more than an order of magnitude improvement in execution speed.
More recently, \texttt{HEAL-SWIN} \cite{carlsson2024healswin} has pushed the state of the art further by adapting the Swin Vision Transformer to the HEALPix sphere: shifted-window attention operates within local HEALPix patches, and the window boundaries are shifted between successive transformer blocks so that the model captures both local and global spherical structure while retaining the computational efficiency of window-based attention.

\paragraph{Cost functions and uncertainty estimation.}
\aure{stopped here}

The choice of cost function determines what the network learns.
For a standard regression problem, the mean squared error $\mathcal{J} = \sum_i (y_i^\mathrm{true} - y_i^\mathrm{pred})^2$ penalizes deviations between predictions and targets uniformly.
When the noise level varies across the output space, however, a more informative approach is to let the network predict both the mean and the variance of its output.
The heteroscedastic Gaussian negative log-likelihood,
\begin{equation}
	\mathrm{NLL} = \sum_{i=1}^{N} \left[ \frac{(y_i^\mathrm{true} - y_i^\mathrm{pred})^2}{2\sigma_i^2} + \frac{1}{2}\ln \sigma_i^2 \right],
	\label{eq:heteroscedastic_nll}
\end{equation}
achieves this through a self-regulating mechanism: the first term allows the network to attenuate the loss for noisy data points by predicting a large $\sigma_i^2$, while the second term penalizes excessive uncertainty, preventing the network from trivially ignoring the data \cite{heteroscedastic-error}.
The variance $\sigma_i^2$ learned in this way captures the \textit{aleatoric} uncertainty---the irreducible noise inherent in the observations, which cannot be reduced by collecting more training data \cite{heteroscedastic-error}.

A second source of uncertainty, called \textit{epistemic} uncertainty, arises from the finite size of the training set and reflects the model's ignorance about which parameter values best describe the data.
Gal and Ghahramani \cite{mcdropout} showed that a neural network with dropout applied before every weight layer is mathematically equivalent to approximate variational inference in a deep Gaussian process.
At test time, this equivalence is exploited through \textit{Monte Carlo dropout}: by performing $T$ stochastic forward passes with dropout active, the sample mean provides the predictive estimate and the sample variance quantifies the epistemic uncertainty \cite{mcdropout}.

\paragraph{Regularization: concrete dropout.}

In standard dropout, the drop probability $p$ must be set manually or tuned via grid search---a procedure that becomes impractical for large architectures.
Concrete dropout \cite{concrete-dropout} addresses this by replacing the discrete Bernoulli mask with a continuous relaxation (the ``concrete'' distribution), making the dropout probability differentiable with respect to the network parameters.
The optimal $p$ for each layer is then learned jointly with the network weights during training, through standard gradient-based optimization.
This self-tuning mechanism simultaneously regularizes the network and calibrates its epistemic uncertainties, without introducing additional hyperparameters.

\paragraph{Application preview.}

In Chapter~\ref{ch:6}, a CNN is applied to gamma-ray sky maps using the heteroscedastic loss and Monte Carlo dropout framework described above, producing both the reconstructed source-count distribution and its associated aleatoric and epistemic uncertainties.
The uncertainty estimates are cross-validated against an independent frequentist analysis (Section~\ref{sec:3.1.4}).

\subsection{Density Estimation}
\label{sec:3.3.3}

The density estimation task introduced in Section~\ref{sec:3.3.1} addresses a complementary problem to classification.
Given a set of $N$ observed data points $\{\mathbf{x}_1, \ldots, \mathbf{x}_N\}$ drawn from an unknown distribution, the goal is to estimate the probability density function $p(\mathbf{x})$ that generated them.
Unlike classification, which learns a decision boundary in feature space, density estimation reconstructs the full probability landscape and can therefore reveal populations that were absent from the training set.

\paragraph{Kernel Density Estimation.}

A non-parametric approach to this problem is provided by Kernel Density Estimation (KDE).
Given $N$ samples, the KDE approximation to the density at a point $\mathbf{x}$ is
\begin{equation}
	\hat{p}(\mathbf{x}) = \frac{1}{N} \sum_{i=1}^{N} K_h(\mathbf{x} - \mathbf{x}_i),
	\label{eq:kde}
\end{equation}
where $K_h$ is a kernel function with bandwidth $h$.
For a Gaussian kernel, $K_h(\mathbf{u}) \propto \exp(-\|\mathbf{u}\|^2 / 2h^2)$, and the bandwidth controls the trade-off between bias and variance: too small a bandwidth overfits the sample, while too large a bandwidth oversmooths the underlying structure.
The optimal bandwidth is determined through cross-validation, selecting the value that maximizes the model likelihood on held-out data.
Separate KDE models are computed independently for each source class (Galactic and extragalactic), yielding the class-conditional densities $p_\mathrm{assoc}(\mathbf{x}|k)$ that form the building blocks of the mixture model described below.

\paragraph{Mixture models and the EM algorithm.}

When the observed data arise from multiple overlapping populations, the total density can be written as a mixture model,
\begin{equation}
	p(\mathbf{x}) = \sum_{k=1}^{K} \pi_k \, p(\mathbf{x}|k),
	\label{eq:mixture_model}
\end{equation}
where $\pi_k$ are the mixing weights (class prevalences) satisfying $\sum_k \pi_k = 1$, and $p(\mathbf{x}|k)$ is the class-conditional density.
The mixing weights and any free parameters of the component densities are typically determined by the Expectation-Maximization (EM) algorithm \cite{Dempster:1977}: an iterative procedure that alternates between computing the posterior class responsibilities $\gamma_{ik} = p(k | \mathbf{x}_i)$ for each data point (the E-step), and updating the model parameters to maximize the expected log-likelihood under those responsibilities (the M-step).
For the mixing weights, the M-step admits a closed-form solution: $\pi_k^\mathrm{new} = N^{-1}\sum_i \gamma_{ik}$, i.e., each weight is updated as the average posterior responsibility over the dataset.

\paragraph{Application preview.}

In Chapter~\ref{ch:5}, this framework is extended to search for dark matter subhalos among the unassociated Fermi-LAT sources.
The mixture model combines astrophysical class-conditional densities with a covariate shift correction (Section~\ref{sec:3.4}) and a possible dark matter component, with class prevalences determined via the EM algorithm following the approach of Saerens et al.~\cite{10.1162/089976602753284446_Saerens}.
